\ifdefined\gkver\else\def\gkver{student}\fi
\documentclass[\gkver]{gaokaozhenti}
\begin{document}

\chapter{2001年北京春季理综}

\begin{enumerate}
\item
%% number: 1
%% paperName: 2001年北京春季理综第 1 题 4 分
%% typeId: 1
%% type: 单选题
%% chapter: 力物体的平衡
%% point: 物体系平衡
%% method: 无
%% score: 4
%% degree: 750
%% duplicateId: 0
%% body:
如图所示，两根相同的轻弹簧 $S_{1}$、$S_{2}$，劲度系数皆为 $k=4\times10^{2}\UNm$。悬挂的重物的质量分别为 $m_{1}=2\Ukg$ 和 $m_{2}=4\Ukg$。若不计弹簧质量，取 $g=10\Umsq$，则平衡时弹簧 $S_{1}$、$S_{2}$ 的伸长量分别为\xzanswer{C}

\onepicture[align=right, w=1.9cm]{figs/01.png}

\fourchoices[answer=C]
{$5\Ucm$、$10\Ucm$}
{$10\Ucm$、$5\Ucm$}
{$15\Ucm$、$10\Ucm$}
{$10\Ucm$、$15\Ucm$}

%% answer: C
%% memo:
\memoanswer{
本题原卷及材料中未提供详解，待补充。
}

\item
%% number: 2
%% paperName: 2001年北京春季理综第 2 题 4 分
%% typeId: 2
%% type: 多选题
%% chapter: 原子和原子核
%% point: 天然放射性
%% method: 无
%% score: 4
%% degree: 850
%% duplicateId: 0
%% body:
下列现象中，与原子核内部变化有关的是\xzanswer{C}

\fourchoices[answer=C]
{$\alpha$ 粒子散射}
{光电效应}
{天然放射现象}
{原子发光现象}

%% answer: C
%% memo:
\memoanswer{
本题原卷及材料中未提供详解，待补充。
}

\item
%% number: 3
%% paperName: 2001年北京春季理综第 3 题 4 分
%% typeId: 1
%% type: 单选题
%% chapter: 气体性质
%% point: 物体的内能
%% method: 无
%% score: 4
%% degree: 650
%% duplicateId: 0
%% body:
下列说法中正确的是\xzanswer{B}

\fourchoices[answer=B]
{物体的分子热运动动能的总和就是物体的内能}
{对于同一种气体，温度越高，分子平均动能越大}
{要使气体的分子平均动能增大，外界必须向气体传热}
{一定质量的气体，温度升高时，分子间的平均距离一定增大}

%% answer: B
%% memo:
\memoanswer{
本题原卷及材料中未提供详解，待补充。
}

\item
%% number: 4
%% paperName: 2001年北京春季理综第 4 题 4 分
%% typeId: 1
%% type: 单选题
%% chapter: 磁场
%% point: 洛伦兹力
%% method: 无
%% score: 4
%% degree: 750
%% duplicateId: 0
%% body:
初速为 $v_{0}$ 的电子，沿平行于通电长直导线的方向射出，直导线中电流方向与电子的初始运动方向如图所示，则\xzanswer{A}

\onepicture[align=right, w=4.2cm]{\includesvg{figs/04}}

\fourchoices[answer=A]
{电子将向右偏转，速率不变}
{电子将向左偏转，速率改变}
{电子将向左偏转，速率不变}
{电子将向右偏转，速率改变}

%% answer: A
%% memo:
\memoanswer{
本题原卷及材料中未提供详解，待补充。
}

\item
%% number: 5
%% paperName: 2001年北京春季理综第 5 题 4 分
%% typeId: 2
%% type: 多选题
%% chapter: 气体性质
%% point: 气体图象
%% method: 无
%% score: 4
%% degree: 650
%% duplicateId: 0
%% body:
一定质量的理想气体经过一系列过程，如图所示。下列说法中正确的是\xzanswer{AD}

\onepicture[align=right, w=4.2cm]{figs/05.png}

\fourchoices[answer=AD]
{$a\to b$ 过程中，气体体积增大，压强减小}
{$b\to c$ 过程中，气体压强不变，体积增大}
{$c\to a$ 过程中，气体压强增大，体积变小}
{$c\to a$ 过程中，气体内能增大，体积不变}

%% answer: AD
%% memo:
\memoanswer{
本题原卷及材料中未提供详解，待补充。
}

\item
%% number: 6
%% paperName: 2001年北京春季理综第 6 题 4 分
%% typeId: 2
%% type: 多选题
%% chapter: 功和能
%% point: 功能中的图象
%% method: 无
%% score: 4
%% degree: 650
%% duplicateId: 0
%% body:
将物体以一定的初速度竖直上抛。若不计空气阻力，从抛出到落回原地的整个过程中，下列四个图线中正确的是\xzanswer{BC}

\fourchoices[answer=BC, ispicture=true, h=2.5cm]
{figs/06a.png}
{figs/06b.png}
{figs/06c.png}
{figs/06d.png}

%% answer: BC
%% memo:
\memoanswer{
本题原卷及材料中未提供详解，待补充。
}

\item
%% number: 7
%% paperName: 2001年北京春季理综第 7 题 4 分
%% typeId: 2
%% type: 多选题
%% chapter: 电路
%% point: 电容
%% method: 无
%% score: 4
%% degree: 650
%% duplicateId: 0
%% body:
一平行板电容器，两板之间的距离 $d$ 和两板面积 $S$ 都可以调节，电容器两板与电池相连接。以 $Q$ 表示电容器的电量，$E$ 表示两极板间的电场强度，则\xzanswer{AC}

\fourchoices[answer=AC]
{当 $d$ 增大、$S$ 不变时，$Q$ 减小、$E$ 减小}
{当 $S$ 增大、$d$ 不变时，$Q$ 增大、$E$ 增大}
{当 $d$ 减小、$S$ 增大时，$Q$ 增大、$E$ 增大}
{当 $S$ 减小、$d$ 减小时，$Q$ 不变、$E$ 不变}

%% answer: AC
%% memo:
\memoanswer{
本题原卷及材料中未提供详解，待补充。
}

\item
%% number: 8
%% paperName: 2001年北京春季理综第 8 题 4 分
%% typeId: 2
%% type: 多选题
%% chapter: 电路
%% point: 动态分析
%% method: 无
%% score: 4
%% degree: 550
%% duplicateId: 0
%% body:
在如图所示的电路中，电容器 $C$ 的上极板带正电。为了使该极板仍带正电且电量增大，下列办法中可采用的是\xzanswer{AD}

\onepicture[align=right, w=4.5cm]{figs/08.png}

\fourchoices[answer=AD]
{增大 $R_{1}$，其他电阻不变}
{增大 $R_{2}$，其他电阻不变}
{增大 $R_{3}$，其他电阻不变}
{增大 $R_{4}$，其他电阻不变}

%% answer: AD
%% memo:
\memoanswer{
本题原卷及材料中未提供详解，待补充。
}

\item
%% number: 9
%% paperName: 2001年北京春季理综第 9 题 4 分
%% typeId: 2
%% type: 多选题
%% chapter: 机械波
%% point: 波的多解性
%% method: 无
%% score: 4
%% degree: 450
%% duplicateId: 0
%% body:
有一列沿水平绳传播的简谐横波，频率为 $10\UHz$，振动方向沿竖直方向。当绳上的质点 P 到达其平衡位置且向下运动时，在其右方相距 $0.6\Um$ 处的质点 Q 刚好到达最高点。由此可知波速和传播方向可能是\xzanswer{BC}

\fourchoices[answer=BC]
{$8\Ums$，向右传播}
{$8\Ums$，向左传播}
{$24\Ums$，向右传播}
{$24\Ums$，向左传播}

%% answer: BC
%% memo:
\memoanswer{
本题原卷及材料中未提供详解，待补充。
}

\item
%% number: 10
%% paperName: 2001年北京春季理综第 10 题 4 分
%% typeId: 2
%% type: 多选题
%% chapter: 运动和力
%% point: 超重失重
%% method: 无
%% score: 4
%% degree: 550
%% duplicateId: 0
%% body:
一物体放置在倾角为 $\theta$ 的斜面上，斜面固定于加速上升的电梯中，加速度为 $a$，如图所示。在物体始终相对于斜面静止的条件下，下列说法中正确的是\xzanswer{BC}

\onepicture[align=right, w=4.5cm]{\includesvg{figs/10}}

\fourchoices[answer=BC]
{当 $\theta$ 一定时，$a$ 越大，斜面对物体的正压力越小}
{当 $\theta$ 一定时，$a$ 越大，斜面对物体的摩擦力越大}
{当 $a$ 一定时，$\theta$ 越大，斜面对物体的正压力越小}
{当 $a$ 一定时，$\theta$ 越大，斜面对物体的摩擦力越小}

%% answer: BC
%% memo:
\memoanswer{
本题原卷及材料中未提供详解，待补充。
}

\item
%% number: 11
%% paperName: 2001年北京春季理综第 11 题 5 分
%% typeId: 3
%% type: 填空题
%% chapter: 原子和原子核
%% point: 人工核反应
%% method: 无
%% score: 5
%% degree: 850
%% duplicateId: 0
%% body:
平衡下列核反应方程式：
\begin{enumerate}
	\item
	\ce{^{y}_{x}B + \alpha{} -> ^{13}_{7}N + n}，$x=$ \tkanswer[1.5cm]{}，$y=$ \tkanswer[2.2cm]{}。
	\item
	\ce{^{7}_{3}Li + p -> \alpha{}} + \tkanswer[2.6cm]{}
\end{enumerate}

%% answer:
\jdanswer{
\begin{enumerate}
	\item
	5，10
	\item
	$\ce{^{4}_{2}He}$ 或 $\alpha$
\end{enumerate}
}
%% memo:
\memoanswer{
本题原卷及材料中未提供详解，待补充。
}

\item
%% number: 12
%% paperName: 2001年北京春季理综第 12 题 5 分
%% typeId: 1
%% type: 单选题
%% chapter: 电场
%% point: 带电粒子的偏转
%% method: 无
%% score: 5
%% degree: 450
%% duplicateId: 0
%% body:
一质量为 $4.0\times10^{-15}\Ukg$、电量为 $2.0\times10^{-9}\UC$ 的带正电质点，以 $4.0\times10^{4}\Ums$ 的速度垂直于电场方向从 a 点进入匀强电场区域，并从 b 点离开电场区域。离开电场时的速度为 $5.0\times10^{4}\Ums$。由此可知，电场中 a、b 两点间的电势差 $U_{a}-U_{b}=$ \tkanswer[3cm]{$9.0\times10^{2}$}$\UV$；带电质点离开电场时，速度在电场方向的分量为 \tkanswer[3cm]{$3.0\times10^{4}$}$\Ums$。不考虑重力作用。

%% answer:
%% memo:
\memoanswer{
本题原卷及材料中未提供详解，待补充。
}

\item
%% number: 13
%% paperName: 2001年北京春季理综第 13 题 5 分
%% typeId: 3
%% type: 填空题
%% chapter: 曲线运动
%% point: 平抛运动
%% method: 无
%% score: 5
%% degree: 450
%% duplicateId: 0
%% body:
质量为 $m=0.10\Ukg$ 的小钢球以 $v_{0}=10\Ums$ 的水平速度抛出，下落 $h=5.0\Um$ 时撞击一钢板，撞后速度恰好反向，则钢板与水平面的夹角 $\theta=$ \tkanswer[2cm]{$45\Udu$}。刚要撞击钢板时小球动量的大小为 \tkanswer[3.4cm]{$\sqrt{2}\Ukgms$ 或 $\sqrt{2}\UNs$}。（取 $g=10\Umsq$）

\onepicture[align=right, w=5.5cm]{figs/13.png}

%% answer:
%% memo:
\memoanswer{
本题原卷及材料中未提供详解，待补充。
}

\item
%% number: 14
%% paperName: 2001年北京春季理综第 14 题 6 分
%% typeId: 4
%% type: 实验题
%% chapter: 直线运动
%% point: 打点计时器公式
%% method: 无
%% score: 6
%% degree: 550
%% duplicateId: 0
%% body:
已知打点计时器接的交流电源频率是 $f$，用它记录一个匀变速运动小车的位移，打出的一条纸带和已选好的计数点 0，1、2、3、4 如图所示。某同学测量出 1 与 2 两点间的距离为 $s_{12}$，3 与 4 两点间的距离为 $s_{34}$，由此可算出小车运动的加速度 $a=$ \tkanswer[5cm]{$\frac{f^{2}}{32}(s_{34}-s_{12})$}。

\onepicture[w=11cm]{figs/14.png}

%% answer:
%% memo:
\memoanswer{
本题原卷及材料中未提供详解，待补充。
}

\item
%% number: 15
%% paperName: 2001年北京春季理综第 15 题 6 分
%% typeId: 4
%% type: 实验题
%% chapter: 光学
%% point: 透镜
%% method: 无
%% score: 6
%% degree: 550
%% duplicateId: 0
%% body:
做测定凸透镜焦距的实验时，把蜡烛和光屏放在透镜的主光轴上，与主光轴垂直。若这时在它们之间无论怎样移动透镜，光屏上都得不到清晰的蜡烛像，则应采取的措施是 \tkanswer[5cm]{加大蜡烛和光屏之间的距离}。为了求得凸透镜的焦距，测出蜡烛到光屏的距离 $L$ 和蜡烛在光屏上两次成像时透镜的两个位置之间的距离 $d$，则该透镜的焦距 $f=$ \tkanswer[3.5cm]{$\frac{L^{2}-d^{2}}{4L}$}。

%% answer:
%% memo:
\memoanswer{
本题原卷及材料中未提供详解，待补充。
}

\item
%% number: 16
%% paperName: 2001年北京春季理综第 16 题 8 分
%% typeId: 4
%% type: 实验题
%% chapter: 电路
%% point: 电阻定律
%% method: 无
%% score: 8
%% degree: 450
%% duplicateId: 0
%% body:
在测定金属电阻率的实验中，用伏安法测量电阻。将实验电路图画在方框中。为了完成整个实验，除你在电路中已画出的器材外，还需要用哪些仪器？用这些仪器测量哪些量？

答：\tkanswer[8cm]{用螺旋测微器测量金属丝的直径 $d$，用米尺测金属丝的长度 $l$}。

计算电阻率 $\rho$ 的公式是

$\rho=$ \tkanswer[3.5cm]{$\frac{\pi U d^{2}}{4Il}$}，式中 $U$ 是电压表测得的电压值，$I$ 是电流表测得的电流值。

%% answer:
\jdanswer{
实验电路图如图所示。

\onepicture[w=4.5cm]{figs/16_an.png}
}
%% memo:
\memoanswer{
本题原卷及材料中未提供详解，待补充。
}

\item
%% number: 17
%% paperName: 2001年北京春季理综第 17 题 11 分
%% typeId: 6
%% type: 计算题
%% chapter: 光学
%% point: 全反射
%% method: 无
%% score: 11
%% degree: 750
%% duplicateId: 0
%% body:
简述光的全反射现象及临界角的定义，并导出折射率为 $n$ 的玻璃对真空的临界角公式。

%% answer:
\jdanswer{
光线从光密介质射向光疏介质时，折射角大于入射角。若入射角增大到某一角度 $C$，使折射角达到 $90\Udu$，折射光就消失。入射角大于 $C$ 时只有反射光，这种现象称为全反射，相应的入射角 $C$ 叫做临界角。

$\sin C=\frac{1}{n}$
}
%% memo:
\memoanswer{
光线从光密介质射向光疏介质时，折射角大于入射角。若入射角增大到某一角度 $C$，使折射角达到 $90\Udu$，折射光就消失。入射角大于 $C$ 时只有反射光，这种现象称为全反射，相应的入射角 $C$ 叫做临界角。

光线由折射率为 $n$ 的玻璃到真空，折射定律为

$\sin r=n\sin i$   ①

其中 $i$ 和 $r$ 分别为入射角和折射角。当入射角 $i$ 等于临界角 $C$ 时，折射角 $r$ 等于 $90\Udu$，代入①式得

$\sin C=\frac{1}{n}$   ②

评分标准：本题 11 分。能正确叙述全反射现象及临界角定义的，给 6 分（如果没有说明全反射，只说临界角是发生全反射的最小入射角的，只给 2 分）。导出②式或 $C=\arcsin\frac{1}{n}$ 的，给 5 分。
}

\item
%% number: 18
%% paperName: 2001年北京春季理综第 18 题 12 分
%% typeId: 6
%% type: 计算题
%% chapter: 曲线运动
%% point: 双星问题
%% method: 无
%% score: 12
%% degree: 550
%% duplicateId: 0
%% body:
两个星球组成双星，它们在相互之间的万有引力作用下，绕连线上某点做周期相同的匀速圆周运动。现测得两星中心距离为 $R$，其运动周期为 $T$，求两星的总质量。

%% answer:
\jdanswer{
$M_{1}+M_{2}=\frac{4\pi^{2}R^{3}}{GT^{2}}$
}
%% memo:
\memoanswer{
设两星质量分别为 $M_{1}$ 和 $M_{2}$，都绕连线上 $O$ 点作周期为 $T$ 的圆周运动，星球 1 和星球 2 到 $O$ 的距离分别为 $l_{1}$ 和 $l_{2}$。由万有引力定律和牛顿第二定律及几何条件可得

$G\frac{M_{1}M_{2}}{R^{2}}=M_{1}(\frac{2\pi}{T})^{2}l_{1}$   ①

$G\frac{M_{1}M_{2}}{R^{2}}=M_{2}(\frac{2\pi}{T})^{2}l_{2}$   ②

$l_{1}+l_{2}=R$   ③

联立解得

$M_{1}+M_{2}=\frac{4\pi^{2}R^{3}}{GT^{2}}$   ④

评分标准：本题 12 分。①、②、③、④各占 3 分。
}

\item
%% number: 19
%% paperName: 2001年北京春季理综第 19 题 12 分
%% typeId: 6
%% type: 计算题
%% chapter: 电路
%% point: 串并联电路
%% method: 无
%% score: 12
%% degree: 450
%% duplicateId: 0
%% body:
如图所示，AB、CD 为两根平行的相同的均匀电阻丝，EF 为另一根电阻丝，其电阻为 $R$，它可以在 AB、CD 上滑动并保持与 AB 垂直，EF 与 AB、CD 接触良好。图中电压表为理想电压表。电池的电动势和内阻都不变。B、D 与电池两极连接的导线的电阻可忽略。当 EF 处于图中位置时，电压表的读数为 $U_{1}=4.0\UV$。已知将 EF 由图中位置向左移动一段距离 $\Delta L$ 后，电压表的读数变为 $U_{2}=3.0\UV$。若将 EF 由图中位置向右移动一段距离 $\Delta L$，电压表的读数 $U_{3}$ 是多少？

\onepicture[align=right, w=6.5cm]{figs/19.png}

%% answer:
\jdanswer{
$U_{3}=6\UV$
}
%% memo:
\memoanswer{
令 $\rho$ 表示 AB、CD 上单位长度电阻丝的电阻，当 EF 处于图中位置时，设 EB、FD 两段电阻丝的长度皆为 $l$。由于电压表是理想电压表，故电压表的读数就是 EF 两端的电压。由分压关系得

$\frac{R\varepsilon}{r+R+2\rho l}=U_{1}$   ①

当 EF 由图示位置向左移动一段距离 $\Delta L$，EB、FD 两段电阻丝的总电阻为 $2\rho(l+\Delta L)$。由分压关系得

$\frac{R\varepsilon}{r+R+2\rho(l+\Delta L)}=U_{2}$   ②

当 EF 由图示位置向右移动一段距离 $\Delta L$，EB、FD 两段电阻丝的总电阻为 $2\rho(l-\Delta L)$。由分压关系得

$\frac{R\varepsilon}{r+R+2\rho(l-\Delta L)}=U_{3}$   ③

由以上三式，代入数值解得

$U_{3}=6\UV$   ④

评分标准：本题 12 分。①、②、③、④式各 3 分。
}

\item
%% number: 20
%% paperName: 2001年北京春季理综第 20 题 12 分
%% typeId: 6
%% type: 计算题
%% chapter: 电磁感应
%% point: 双杆问题
%% method: 无
%% score: 12
%% degree: 350
%% duplicateId: 0
%% body:
两根足够长的固定的平行金属导轨位于同一水平面内，两导轨间的距离为 $l$。导轨上面横放着两根导体棒 ab 和 cd，构成矩形回路，如图所示。两根导体棒的质量皆为 $m$，电阻皆为 $R$，回路中其余部分的电阻可不计。在整个导轨平面内都有竖直向上的匀强磁场，磁感应强度为 $B$。设两导体棒均可沿导轨无摩擦地滑行。开始时，棒 cd 静止，棒 ab 有指向棒 cd 的初速度 $v_{0}$（见图）。若两导体棒在运动中始终不接触，求：
\begin{enumerate}
	\item
	在运动中产生的焦耳热最多是多少；
	\item
	当 ab 棒的速度变为初速度的 $\frac{3}{4}$ 时，cd 棒的加速度是多少？
\end{enumerate}

\twopicture[align=right, w=6.0cm, labelA=2001北京春季理综20a, labelB=2001北京春季理综20b]
{figs/20a.jpg}
{figs/20b.png}

%% answer:
\jdanswer{
\begin{enumerate}
	\item
	$Q=\frac{1}{4}mv_{0}^{2}$
	\item
	$a=\frac{B^{2}l^{2}v_{0}}{4mR}$
\end{enumerate}
}
%% memo:
\memoanswer{
ab 棒向 cd 棒运动时，两棒和导轨构成的回路面积变小，磁通量发生变化，于是产生感应电流。ab 棒受到与运动方向相反的安培力作用作减速运动，cd 棒则在安培力作用下作加速运动。在 ab 棒的速度大于 cd 棒的速度时，回路总有感应电流，ab 棒继续减速，cd 棒继续加速。两棒速度达到相同后，回路面积保持不变，磁通量不变化，不产生感应电流，两棒以相同的速度 $v$ 作匀速运动。

（1）从初始至两棒达到速度相同的过程中，两棒总动量守恒，有

$mv_{0}=2mv$   ①

根据能量守恒，整个过程中产生的总热量

$Q=\frac{1}{2}mv_{0}^{2}-\frac{1}{2}(2m)v^{2}=\frac{1}{4}mv_{0}^{2}$   ②

（2）设 ab 棒的速度变为初速度的 $\frac{3}{4}$ 时，cd 棒的速度为 $v'$，则由动量守恒可知

$mv_{0}=m\cdot\frac{3}{4}v_{0}+mv'$   ③

此时回路中的感应电动势和感应电流分别为

$\varepsilon=(\frac{3}{4}v_{0}-v')Bl$   ④

$I=\frac{\varepsilon}{2R}$   ⑤

此时 cd 棒所受的安培力

$F=IBl$   ⑥

cd 棒的加速度

$a=\frac{F}{m}$   ⑦

由以上各式，可得

$a=\frac{B^{2}l^{2}v_{0}}{4mR}$   ⑧

评分标准：本题 12 分。第（1）问 6 分，其中①、②各 3 分。第（2）问 6 分，其中③式 1 分，④式 2 分，⑤式 1 分，⑧式 2 分。
}

\item
%% number: 21
%% paperName: 2001年北京春季理综第 21 题 14 分
%% typeId: 6
%% type: 计算题
%% chapter: 气体性质
%% point: 气体多过程
%% method: 无
%% score: 14
%% degree: 350
%% duplicateId: 0
%% body:
如图所示，一水平放置的气缸，由截面积不同的两圆筒联接而成。活塞 A、B 用一长为 $3l$ 的刚性细杆连接，它们可以在筒内无摩擦地沿水平方向左右滑动。A、B 的截面积分别为 $S_{A}=30\Ucmq$、$S_{B}=15\Ucmq$。A、B 之间封闭着一定质量的理想气体。两活塞外侧（A 的左方和 B 的右方）都是大气，大气压强始终保持为 $p_{0}=1.0\times10^{5}\UPa$。活塞 B 的中心连一不能伸长的细线，细线的另一端固定在墙上。当气缸内气体温度为 $T_{1}=540\UK$，活塞 A、B 的平衡位置如图所示，此时细线中的张力为 $F_{1}=30\UN$。
\begin{enumerate}
	\item
	现使气缸内气体温度由初始的 $540\UK$ 缓慢下降，温度降为多少时活塞开始向右移动？
	\item
	继续使气缸内气体温度下降，温度降为多少时活塞 A 刚刚右移到两圆筒联接处？
	\item
	活塞 A 移到两圆筒联接处之后，维持气体温度不变，另外对 B 施加一个水平向左的推力，将两活塞慢慢推向左方，直到细线拉力重新变为 $30\UN$。求此时的外加推力 $F_{2}$ 是多大。
\end{enumerate}

\onepicture[align=right, w=7.5cm]{\includesvg{figs/21}}

%% answer:
\jdanswer{
\begin{enumerate}
	\item
	$T_{2}=450\UK$
	\item
	$T_{3}=270\UK$
	\item
	$F_{2}=90\UN$
\end{enumerate}
}
%% memo:
\memoanswer{
（1）设气缸内气体压强为 $p_{1}$，$F$ 为细线中的张力，则活塞 A、B 及细杆这个整体的平衡条件为

$p_{0}S_{A}-p_{1}S_{A}+p_{1}S_{B}-p_{0}S_{B}+F_{1}=0$（1分），解得 $p_{1}=p_{0}+\frac{F_{1}}{S_{A}-S_{B}}$   ①（1分）

对于初始状态，$F=F_{1}=30\UN$，

代入①式，就得到气缸中气体的初始压强 $p_{1}=1.2\times10^{5}\UPa$（1分）

由①式看出，只要气体压强 $p_{1}>p_{0}$，细线就会拉直且有拉力，于是活塞不会移动。使气缸内气体温度降低是等容降温过程，当温度下降使压强降到 $p_{0}$ 时，细线拉力变为零，再降温时活塞开始向右移，设此时温度为 $T_{2}$，压强 $p_{2}=p_{0}$。有

$\frac{T_{2}}{T_{1}}=\frac{p_{0}}{p_{1}}$   ②（1分）

得 $T_{2}=\frac{5}{6}T_{1}=450\UK$   ③（1分）

（2）再降温，细线松了，要平衡必有气体压强 $p=p_{0}$。是等压降温过程，活塞右移、体积相应减小，当 A 到达两圆筒联接处时，温度为 $T_{3}$，

由 $\frac{V_{2}}{T_{2}}=\frac{V_{3}}{T_{3}}$

即 $\frac{2lS_{A}+lS_{B}}{T_{2}}=\frac{3S_{B}l}{T_{3}}$   ④（2分）

得 $T_{3}=\frac{1}{2}T_{1}=270\UK$   ⑤（1分）

（3）维持 $T_{3}=270\UK$ 不变，向左推活塞，是等温过程，最后压强为 $p_{4}$。有

$\frac{p_{4}}{p_{0}}=\frac{3lS_{B}}{2lS_{A}+lS_{B}}$   ⑥（2分）

推力 $F_{2}$ 向左，由力的平衡条件得

$p_{0}S_{A}-p_{4}S_{A}+p_{4}S_{B}-p_{0}S_{B}+F_{1}-F_{2}=0$   ⑦（1分）

解得 $F_{2}=90\UN$   ⑧（1分）

评分标准：本题 14 分。第（1）问 6 分，其中①、②、③式各 2 分。第（2）问 4 分，其中④、⑤式各 2 分。第（3）问 4 分，其中⑥式 2 分，⑦式 1 分，⑧式 1 分。
}

\item
%% number: 22
%% paperName: 2001年北京春季理综第 22 题 14 分
%% typeId: 6
%% type: 计算题
%% chapter: 运动和力
%% point: 动量和能量
%% method: 无
%% score: 14
%% degree: 250
%% duplicateId: 0
%% body:
如图所示，A、B 是静止在水平地面上完全相同的两块长木板。A 的左端和 B 的右端相接触。两板的质量皆为 $M=2.0\Ukg$，长度皆为 $l=1.0\Um$。C 是一质量为 $m=1.0\Ukg$ 的小物块。现给它一初速度 $v_{0}=2.0\Ums$，使它从 B 板的左端开始向右滑动。已知地面是光滑的，而 C 与 A、B 之间的动摩擦因数皆为 $\mu=0.10$。求最后 A、B、C 各以多大的速度做匀速运动。取重力加速度 $g=10\Umsq$。

\onepicture[align=right, w=6.5cm]{\includesvg{figs/22}}

%% answer:
\jdanswer{
$v_{B}=0.155\Ums$，$v_{C}=v_{A}=0.563\Ums$
}
%% memo:
\memoanswer{
先假设小物块 C 在木板 B 上移动 $x$ 距离后，停在 B 上。这时 A、B、C 三者的速度相等，设为 $V$。由动量守恒得

$mv_{0}=(m+2M)V$   ①

在此过程中，木板 B 的位移为 $s$，小物块 C 的位移为 $s+x$。由功能关系得

$-\mu mg(s+x)=\frac{1}{2}mV^{2}-\frac{1}{2}mv_{0}^{2}$

$\mu mgs=\frac{1}{2}\cdot2MV^{2}$

相加得

$-\mu mgx=\frac{1}{2}(m+2M)V^{2}-\frac{1}{2}mv_{0}^{2}$   ②

解①、②两式得

$x=\frac{Mv_{0}^{2}}{(2M+m)\mu g}$   ③

代入数值得

$x=1.6\Um$   ④

$x$ 比 B 板的长度 $l$ 大。这说明小物块 C 不会停在 B 板上，而要滑到 A 板上。设 C 刚滑到 A 板上的速度为 $v_{1}$，此时 A、B 板的速度为 $V_{1}$，则由动量守恒得

$mv_{0}=mv_{1}+2MV_{1}$   ⑤

由功能关系得

$\frac{1}{2}mv_{0}^{2}-\frac{1}{2}mv_{1}^{2}-\frac{1}{2}\cdot2MV_{1}^{2}=\mu mgl$   ⑥

以题给数据代入解得

$V_{1}=\frac{8\pm\sqrt{24}}{20}$

$v_{1}=2-\frac{8\pm\sqrt{24}}{5}=\frac{2\mp\sqrt{24}}{5}$

由于 $v_{1}$ 必是正数，故合理的解是

$V_{1}=\frac{8-\sqrt{24}}{20}=0.155\Ums$   ⑦

$v_{1}=\frac{2+\sqrt{24}}{5}=1.38\Ums$   ⑧

当滑到 A 之后，B 即以 $V_{1}=0.155\Ums$ 做匀速运动。而 C 是以 $v_{1}=1.38\Ums$ 的初速在 A 上向右运动。设在 A 上移动了 $y$ 距离后停止在 A 上，此时 C 和 A 的速度为 $V_{2}$，由动量守恒得

$MV_{1}+mv_{1}=(m+M)V_{2}$   ⑨

解得

$V_{2}=0.563\Ums$   ⑩

由功能关系得

$\frac{1}{2}mv_{1}^{2}+\frac{1}{2}MV_{1}^{2}-\frac{1}{2}(m+M)V_{2}^{2}=\mu mgy$   \textcircled{\scriptsize 11}

解得 $y=0.50\Um$   \textcircled{\scriptsize 12}

$y$ 比 A 板的长度小，故小物块 C 确实是停在 A 板上。最后 A、B、C 的速度分别为 $v_{A}=v_{2}=0.563\Ums$，$v_{B}=v_{1}=0.155\Ums$，$v_{C}=v_{A}=0.563\Ums$。

评分标准：本题 14 分。正确论证了 C 不能停在 B 板上而是停在 A 板上，占 8 分。求出 A、B、C 三者的最后速度，占 6 分。
}

\end{enumerate}

\end{document}
